% ============================================================
%  style/illustrationen.tex — Illustrationen frei auf der Seite platzieren
%  (wird von style.tex eingebunden)
% ============================================================
% ---------- Illustrationen frei platzieren ----------
% \illustration[Breite]{Datei}{x}{y}[includegraphics-Optionen]
%  - Datei: PDF in graphics/illustrationen/ (ohne .pdf)
%  - x / y: Abstand der linken oberen Bildecke von der linken oberen
%    PAPIERECKE (absolut, z.B. 95mm / 20mm). Achtung: linke/rechte Seiten
%    haben unterschiedliche Ränder.
%  - Breite optional (Standard 40mm), Höhe ergibt sich aus dem Seitenverhältnis
%  - letzte Option z.B. "trim=163 79 156 36, clip" um Weißraum
%    abzuschneiden (Werte in pt: links unten rechts oben)
%  - Liegt IM VORDERGRUND (über Text/Akkorden).
%  - Erscheint auf der Seite, auf der die Zeile mit dem Befehl landet
%    (z.B. hinten an eine Songzeile hängen). Seite wird über die .aux
%    ermittelt -> zwei Durchläufe (latexmk erledigt das).
\makeatletter
\NewDocumentCommand{\illustration}{ O{40mm} m m m O{} }{%
  \edef\illu@tmp{\write\@auxout{\string\illustrationseite
    {\noexpand\the\ReadonlyShipoutCounter}%
    {\detokenize{#1}}{\detokenize{#2}}{\detokenize{#3}}{\detokenize{#4}}%
    {\detokenize{#5}}}}%
  \illu@tmp
}
% aus der .aux: Illustration für Seite #1 vormerken
\newcommand{\illustrationseite}[6]{%
  \csgappto{illu@seite@#1}{\illustrationzeichnen{#2}{#3}{#4}{#5}{#6}}%
}
\newcommand{\illustrationzeichnen}[5]{%
  \AtPageUpperLeft{\put(\LenToUnit{#3},\LenToUnit{-#4}){%
    \raisebox{-\height}{\includegraphics[width=#1,#5]{graphics/illustrationen/#2.pdf}}}}%
}
% beim zweiten Lesen der .aux (\end{document}) nicht doppelt eintragen
\AtEndDocument{\renewcommand{\illustrationseite}[6]{}}
\AddToShipoutPictureFG{\csuse{illu@seite@\the\ReadonlyShipoutCounter}}
\makeatother

